% Generated draft: requires mathematical and source review.
\documentclass[11pt]{article}
\usepackage[margin=25mm]{geometry}
\usepackage{fontspec}
\setmainfont{Latin Modern Roman}
\usepackage{amsmath,amssymb,mathtools,unicode-math}
\setmathfont{Latin Modern Math}
\usepackage{enumitem,needspace,xurl,hyperref}
\setlength{\parindent}{0pt}
\setlength{\parskip}{5pt}
\setlength{\emergencystretch}{3em}
\newcommand{\N}{\mathbb{N}}
\newcommand{\Z}{\mathbb{Z}}
\newcommand{\Q}{\mathbb{Q}}
\newcommand{\R}{\mathbb{R}}
\newcommand{\C}{\mathbb{C}}
\newcommand{\F}{\mathbb{F}}
\newcommand{\A}{\mathbb{A}}
\newcommand{\PP}{\mathbb{P}}
\newcommand{\E}{\mathbb{E}}
\newcommand{\eps}{\varepsilon}
\newcommand{\dd}{\,\mathrm{d}}
\begin{document}
\section*{UnsolvedMath ID: 700754}
\textbf{Source problem code:} catalogue entry
\subsection{Definitions and mathematical statement}
The supplied record does not define ``equation'' or ``soluble over a group.'' Consequently, the following is a conditional formalization using a one-variable equation and a solution in an overgroup; these conventions require confirmation against the cited source.

Let \(G\) be a group, let \(\langle x\rangle\cong\Z\), and form the free product \(G*\langle x\rangle\). For this provisional formulation, an admissible equation over \(G\) is
\[
 w(x)=1,\qquad w\in G*\langle x\rangle,
\]
where \(w\) is not conjugate to an element of the coefficient subgroup \(G\). This restriction excludes equations that are merely conjugates of constant equations; it is an interpretive convention, not an explicit hypothesis in the supplied record. No restriction on the exponent sum of \(x\) is imposed here.

Such an equation is called soluble over \(G\) if there exist a group \(H\), an injective homomorphism \(\iota:G\hookrightarrow H\), and an element \(h\in H\) such that
\[
 \operatorname{ev}_{\iota,h}(w)=1_H,
\]
where \(\operatorname{ev}_{\iota,h}:G*\langle x\rangle\to H\) restricts to \(\iota\) on \(G\) and sends \(x\) to \(h\).

With these provisional conventions, define
\[
 \mathcal C=\left\{G:\text{every admissible equation over }G\text{ is soluble over }G\right\}.
\]
The problem is to characterize \(\mathcal C\). In particular, is it true that, for every group \(G\),
\[
 G\in\mathcal C
 \quad\Longleftrightarrow\quad
 G\text{ is torsion-free}?
\]
Here torsion-free means that, for every \(g\in G\) and every integer \(n\geq1\), \(g^n=1_G\) implies \(g=1_G\).

The supplied statement says that membership in the intended class implies torsion-freeness. It also attributes to S. D. Brodski\u{\i} the assertion that every equation is soluble over a locally indicable group. A group is locally indicable if each of its nontrivial finitely generated subgroups admits a surjective homomorphism onto \(\Z\). These assertions are retained as supplied background; their exact equation conventions have not been verified. The registry labels the problem open, which is supplied curation rather than an independently verified status assessment.
\subsection{Short English statement}
Characterize the groups over which every equation is soluble, and ask whether they are exactly the torsion-free groups. The draft provisionally interprets this as solving admissible one-variable equations in an overgroup containing the coefficient group injectively. The source does not specify these conventions, so the intended scope requires review. It supplies torsion-freeness as necessary and attributes the locally indicable case to S. D. Brodski\u{\i}.
\subsection{Sources}
\begin{itemize}
\item Supplied local UnsolvedMath registry record: permanent ID \(700754\), exact source code \texttt{catalogue entry}, title ``Kourovka Notebook Problem 11.62.'' The record labels the problem \texttt{open}; this label has not been independently verified.
\item \emph{The Kourovka Notebook (Unsolved Problems in Group Theory)}, problem 11.62, originating in the 11th issue (1990); the supplied background identifies the 21st-issue edition (2026), \texttt{arXiv:1401.0300v48}. The supplied collection description names E. I. Khukhro and V. D. Mazurov as editors and covers Issues 1--21 (1965--2026). No external source was consulted.
\item Proposed in 1990 by the person recorded literally as \texttt{D. I. Moldavanski\textbackslash{}u\{\textbackslash{}i}. The supplied proposer string contains an incomplete TeX accent sequence; its spelling has not been silently repaired.
\item S. D. Brodski\u{\i}, \emph{Siberian Math. J.}, volume 25, no. 2 (1984), pp. 235--251. The supplied statement cites this reference for solubility of every equation over a locally indicable group. No article title or URL is supplied, and the precise theorem and its hypotheses have not been independently checked.
\end{itemize}
\end{document}
